\section*{Chapter \ref{ch:g10:images-sound-compression} --- Digital Images, Sound and Compression}

\begin{solution}{exo:g10:images-sound-compression:1}
$4000 \times 3000 \times 3 = 36{,}000{,}000$ bytes, that is $36{,}000{,}000 / 2^{20} \approx 34.3$~MiB, as the listing
printed. In grey levels, one byte per pixel: $12{,}000{,}000$ bytes, about $11.4$~MiB.
\end{solution}

\begin{solution}{exo:g10:images-sound-compression:2}
$48{,}000 \times 2 \times 2 = 192{,}000$ bytes per second. One hour: $3600 \times 192{,}000 = 691{,}200{,}000$ bytes,
about $659$~MiB.
\end{solution}

\begin{solution}{exo:g10:images-sound-compression:3}
\code{[(3, "A"), (1, "B"), (4, "C"), (1, "D")]} (Python prints the letters with single quotes): $4$ numbers and $4$
characters, $8$ values against $9$. Runs of length $1$ cost more than they save.
\end{solution}

\begin{solution}{exo:g10:images-sound-compression:4}
The $10{,}000$ letters \code{a} come out smallest, a few dozen bytes: the whole input is one repetition. The chapter's
text comes next, a few thousand bytes: words and spaces repeat, and some letters are much more frequent than others.
The random bytes come out largest, a little more than $10{,}000$ bytes: there is nothing to exploit, and the format
adds a few bytes of its own, as the listing showed.
\end{solution}

\begin{solution}{exo:g10:images-sound-compression:5}
\code{return "".join(ch * count for count, ch in runs)}: each pair gives its character repeated \code{count} times,
and the pieces are joined in order. The listing checks \code{rle_decode(rle_encode(row)) == row}; the chapter's tests
check it on $300$ random strings.
\end{solution}

\begin{solution}{exo:g10:images-sound-compression:6}
A run is appended only when the character changes, and the last run is followed by no change: it is never appended.
After the loop, add the last run (\code{exercises/rle_bugs.py}):

\smallskip\noindent\hspace*{2em}\code{if s: runs.append((count, s[-1]))}

\smallskip\noindent The test on the empty string keeps an empty list. The buggy version is right only on the empty string: every other string has a last run.
\end{solution}

\begin{solution}{exo:g10:images-sound-compression:7}
The worst row has no two equal neighbours, such as \code{WBWBWB}: $n$ runs of length $1$, so $2n$ bytes, twice the
row. The encoding is shorter exactly when there are fewer than $n/2$ runs, that is, when runs are on average longer
than $2$ pixels: plain areas, drawings and maps, not photographs.
\end{solution}

\begin{solution}{exo:g10:images-sound-compression:8}
Split the text into words; check that the first is \code{P2}; read the width, height and largest value; take the
next \code{width * height} words as whole numbers; count them in a list of \code{maxval + 1} zeros
(\code{exercises/pgm_histogram.py}). For the image \code{P2 3 2 4} with values \code{0 1 1 4 4 4} it returns
\code{[1, 2, 0, 0, 3]}.
\end{solution}

\begin{solution}{exo:g10:images-sound-compression:9}
$2^{16} = 65{,}536$ colours. Red and blue have $2^5 = 32$ levels each, green $2^6 = 64$. The eye tells shades of green
apart best, as the colour filter mosaic also reflects with its two green \omterm{def:g10:images-sound-compression:photosite}{photosites} in four, so the extra bit is best
spent there.
\end{solution}

\begin{solution}{exo:g10:images-sound-compression:10}
With \code{<= 4}, the index $4$ is accepted, but the rows and columns are numbered $0$ to $3$. For every pixel in the
last row or the last column, a neighbour at index $4$ is read, and \code{MOSAIC[4]} or \code{MOSAIC[r][4]} raises an
\code{IndexError}: the program stops. The other pixels are unchanged. The condition must be \code{< 4}, or better
\code{< len(MOSAIC)}.
\end{solution}

\begin{solution}{exo:g10:images-sound-compression:11}
Open the file with \code{wave.open(path, "rb")} and return \code{w.getnframes() / w.getframerate()}
(\code{exercises/wav_duration.py}): a frame holds one sample per channel, so the number of frames per second is the
\omterm{def:g10:images-sound-compression:sampling-rate}{sampling rate}. The tone of the listing has $8000$ frames at $8000$~Hz: $1.0$~s.
\end{solution}

\begin{solution}{exo:g10:images-sound-compression:12}
There are $256^n$ files of exactly $n$ bytes, and $1 + 256 + 256^2 + \dots + 256^{n-1}$ files of fewer than $n$ bytes.
Multiplying this sum by $256 - 1$ gives $256^n - 1$, since all the other terms cancel; so the sum is
$(256^n - 1)/255$, less than $256^n$. A \omterm{def:g10:images-sound-compression:compression}{lossless} method must send two different files to two different results,
otherwise decompression could not tell which one to give back. If it made every file of $n$ bytes shorter, it would
send $256^n$ different files into fewer than $256^n$ shorter files: two of them would share a result. So some file of
$n$ bytes is not made shorter.
\end{solution}

\begin{solution}{exo:g10:images-sound-compression:13}
\code{bitmap_ratio} adds the lengths of the rows and the numbers of runs of each row, and returns the pixels divided
by twice the runs (\code{exercises/rle_bitmap.py}). On a map with large plain areas the ratio is well above $1$. On an
$8 \times 8$ chessboard every pixel is a run: $64$ pixels, $64$ runs, a ratio of $0.5$, the worst case of
\cref{exo:g10:images-sound-compression:7}.
\end{solution}

\begin{solution}{exo:g10:images-sound-compression:14}
The levels are $0, s, 2s, \dots$. A grey value $v$ lies between two neighbouring levels $ks$ and $(k + 1)s$, so its
distances to them add up to $s$; the smaller of the two is at most $s/2$, and rounding picks the nearer level. In the
listing the levels are also rounded to whole numbers, which can add up to $\tfrac12$. For $16$ levels, $s = 17$ and
the largest error printed is $8 \le 8.5$; for $8$ levels, $s \approx 36.43$ and the largest error is
$18 \le 18.2$.
\end{solution}

\begin{solution}{exo:g10:images-sound-compression:15}
The count gives $1000$ crossings per second (\code{exercises/aliasing.py}): the recording seems to be a $1000$~Hz
tone. Between two samples the $7000$~Hz wave makes $7000/8000 = \tfrac78$ of a turn, which lands at the same place as
going back $\tfrac18$ of a turn. The samples therefore trace a wave that goes back one eighth of a turn per sample,
one turn every $8$ samples: $8000 / 8 = 1000$ turns per second. The samples cannot tell the two waves apart, which is
why the \omterm{def:g10:images-sound-compression:sampling-rate}{sampling rate} must be more than twice the highest frequency kept.
\end{solution}

\begin{solution}{pb:g10:images-sound-compression:1}
\textbf{1.} $44{,}100$ samples per second, $16$ bits per sample.

\textbf{2.} $2^{16} = 65{,}536$ levels.

\textbf{3.} $2 \times 44{,}100 = 88{,}200$ samples.

\textbf{4.} $1/44{,}100$~s $\approx 22.7$~µs.

\textbf{5.} $88{,}200 \times 2 = 176{,}400$ bytes per second, as the program prints.

\textbf{6.} $60 \times 176{,}400 = 10{,}584{,}000$ bytes.

\textbf{7.} $10{,}584{,}000 / 2^{20} \approx 10.09$~MiB, and $10.584$~MB.

\textbf{8.} $8 \times 176{,}400 = 1{,}411{,}200$ bits per second.

\textbf{9.} $210 \times 176{,}400 = 37{,}044{,}000$ bytes, about $35.3$~MiB.

\textbf{10.} $10^9 / 37{,}044{,}000 \approx 26.99$: $26$ whole songs.

\textbf{11.} $44 / 37{,}044{,}000 \approx 1.2 \times 10^{-6}$, about a millionth: nothing.

\textbf{12.} $1{,}411{,}200 / 160{,}000 = 8.82$.

\textbf{13.} \omterm{def:g10:images-sound-compression:compression}{Lossy}. Part IV measures a \omterm{def:g10:images-sound-compression:compression}{lossless} ratio of about $1.09$ on this music with \code{zlib}, nowhere near $8.82$. Methods built
for sound do better by predicting each sample from the ones before, but still nowhere near $8.82$; and the
proposition says that no \omterm{def:g10:images-sound-compression:compression}{lossless} method shrinks every file. The service throws away what the ear hardly notices.

\textbf{14.} $175{,}662$ runs for $176{,}400$ bytes: almost every byte is a run of its own. \omterm{def:g10:images-sound-compression:rle}{Run-length coding} with
two bytes per run would nearly double the size, to about $351{,}000$ bytes.

\textbf{15.} About $1.09$ on music ($161{,}664$ bytes) and about $909$ on silence ($194$ bytes).

\textbf{16.} A sample is stored in two bytes, the low byte first. The low byte holds the finest $\tfrac{1}{256}$ of
the amplitude: even a smooth tone moves by hundreds of levels from one sample to the next, and the small noise
changes it at random. The bytes of the two channels and of the two halves of each sample also alternate. Seen as a
sequence of bytes, the music looks almost random.

\textbf{17.} Every byte is $0$: one run of $176{,}400$ bytes, which a dictionary method describes in a few
references to what it has just written.

\textbf{18.} The silent part compresses as well as the silence of Part IV, so the file shrinks by roughly the length of
that silence; the music before it hardly shrinks at all.

\textbf{19.} One minute of disc-quality stereo is $44{,}100 \times 2 \times 2 \times 60 = 10{,}584{,}000$ bytes, about
$10.1$~MiB. A general-purpose \omterm{def:g10:images-sound-compression:compression}{lossless} compressor, \code{zlib}, reaches a ratio of only about $1.09$ on this music, against
about $909$ on silence: seen as bytes, the samples look nearly random. Compressors made for sound predict each sample and
do better, but not by a factor of eight: that is why music files are made small with \omterm{def:g10:images-sound-compression:compression}{lossy} methods.
\end{solution}
